\section{What the studies support}\label{sec:discussion}

\begin{table}[!htb]\centering
\caption{Standing of the eight dimensions after both studies. For the second study: the reliability of the
composite and the better model's correlation with it.}\label{tab:standing}
\footnotesize
\begin{tabularx}{\textwidth}{@{}l>{\raggedright\arraybackslash}p{2.9cm}>{\raggedright\arraybackslash}p{2.6cm}>{\raggedright\arraybackslash}p{1.6cm}>{\raggedright\arraybackslash}p{3.2cm}>{\raggedright\arraybackslash}X@{}}
\toprule
 & Dimension & First study,\newline by question\newline (fixed in advance) & First study,\newline both markers\newline (added later) & Second study,\newline by dimension\newline (fixed in advance) & Standing \\
\midrule
F1 & Direct, organized argument & both open & usable & \textbf{usable}, settled\newline 0.71; 0.72 & supported \\
F2 & Positive tone versus disappointment & both usable & usable & \textbf{usable}, settled\newline 0.86; 0.79 & supported twice \\
F3 & Self-disclosure and self-criticism & both usable & usable & \textbf{usable}, settled\newline 0.84; 0.81 & supported twice \\
F4 & Care and regard for others & borderline\newline (compassion) & & usable, at the line\newline 0.66; 0.50 & open: followed weakly, and not alike in streams and videos \\
F5 & Narrative versus rebuttal & open\newline (narrative) & & borderline, at the line\newline 0.60; 0.70 & open: the models follow, raters at the line \\
F6 & Energetic and angry versus calm & borderline, open & usable & \textbf{usable}, settled\newline 0.77; 0.82 & supported \\
F7 & Humor and imaginative language & borderline\newline (humor) & & borderline, at the line\newline 0.59; 0.56 & open \\
F8 & Hedging and open-mindedness versus certainty & borderline\newline (hedging) & & unclear rubric\newline 0.34; 0.58 & not shown; one question is not readable \\
\bottomrule
\end{tabularx}
\end{table}

\paragraph{Supported, twice.} On positive tone, disappointment, self-critical wording and self-disclosure, either
model ranks 250-word excerpts much as five readers do. This was fixed in advance in the first study and held
again in the second, on other excerpts, with other raters and in another design. The two dimensions these
questions mark, F2 and F3, are usable and settled (Table~\ref{tab:standing}).

\paragraph{Supported, found first and then confirmed.} The dimension scores of F1 and F6 follow people's ratings
of their two markers taken together. The first study found this after its results had been read; the second
named it in advance as what would count as confirmation, and found both dimensions usable with the whole
interval above the lines.

\paragraph{Supported by the analyses added afterwards, with less weight.} People's ratings show the same
differences between videos and streams, and the same changes over the years, as the models' scores. A model is
closer to what many readers would say together than one reader is. Each dimension score follows people's
composite of its own dimension more closely than that of any other. The models follow the difference between
two passages of the same recording. How well the models follow people on a question repeats from one study
to the next. And a third model from another maker follows people to the same degree and misses where the first
two miss, so the findings are not a property of one model.

\paragraph{Open.} Three dimensions lie at the lines of the rule, each for another reason.
\begin{itemize}
\item \textbf{F5, narrative versus rebuttal.} Both models follow the composite well. What is at the line is the
  raters' agreement, and the two markers hardly go together in people's reading.
\item \textbf{F7, humor and imaginative language.} Humor alone is rated at least as reliably as in the first
  study and followed by the models; imaginative language is not agreed on.
\item \textbf{F4, care and regard for others.} Raters agree well enough. The models' correlation is at its
  line, the dimension score does not clearly follow its own composite more than another, people read the
  dimension as opposed to energetic and angry wording where the models' scores do not, and at the same score
  people rate streams half a standard deviation lower than videos. The comparison of streams with videos is
  not supported for this dimension.
\end{itemize}

\paragraph{Not shown.} F8, hedging and open-mindedness versus certainty. Raters do not agree on its second
marker, knowledge limits, at all, and only weakly on hedging in either study. The models treat the two
questions as nearly one; people do not. The studies do not show that the models misread the dimension: they
show that this pair of questions does not measure it in people.

\paragraph{Not supported.} Reading a model's level as the level a person would give. In both studies one model
rates hedging about a scale point higher than people, and the other rates most questions lower. The models
rank like people where a dimension is usable; their scales are their own.

\subsection{Limits}\label{sec:limits}

\begin{itemize}
\item \textbf{What was rated.} The first 250 words of a 500-word passage of a machine transcript, as plain text.
  The archive is rated at 500 words. Nothing was heard, and nothing marks who speaks. The figures for
  recordings rest on two such passages per recording.
\item \textbf{Which questions.} Sixteen of the 48 rating questions. The other 32, and the classifications, were
  not rated by people.
\item \textbf{Who rated.} Crowd workers of one platform who read English and have a high approval rate, five per
  excerpt and question. In the second study they are people who took part after a warning about swearing and
  sexual language. Nothing else is known about them. The consensus is what this population says on average.
\item \textbf{Which models.} \jev{}, \drex{} and, added afterwards, d1.2-ultra, as they answered in November
  2026, at this length of text. Dimension scores of the two later models go through a scorer fitted on \jev{}.
\item \textbf{``Cannot judge''} is treated as missing (2.0\% and 1.9\% of answers).
\item \textbf{The first study's figures} rest on the assumptions of Section~\ref{sec:firstmethods}, and its
  reading check was changed during collection; both versions are reported.
\item \textbf{The second study's intervals} describe which recordings were drawn; the choice of raters adds
  about as much (Section~\ref{sec:firm}). Its analysis was run once before the last sets were filled, and one
  verdict changed sides.
\item \textbf{The lines of the verdicts} are conventions fixed in advance. Three of the eight dimensions lie
  within their own interval of a line.
\item \textbf{Analyses added afterwards} were not fixed in advance. With many figures some differences will look
  clear by chance.
\item \textbf{What the ratings are.} Ratings of wording by readers who saw nothing but the text. They are not
  measurements of a person, and a passage can hold other voices than the one a reader assumes.
\end{itemize}

\subsection{What stays open}\label{sec:open}

No further text study is planned. What another study of this kind could answer is specific to single
dimensions and single questions, and the general questions are answered: the models rank excerpts as readers do
where readers agree, they follow the passage and not only the recording, their two voices add little to each
other, and their levels are their own. Open are the three dimensions at the line; the three questions that
readers could not rate reliably (knowledge limits, imaginative language, respect), for which other wording
would be needed; why readers and models differ on care and regard for others in livestreams; and everything
about length and sound, since 250 words of plain text were rated and nothing was heard.
